\section{Entire face maps and edge returns}
\label{cg:maps}

The construction matches complete faces isometrically. We now
give the complex-coordinate dictionary for their identifications.
It records the shear and checks the entire edge return, including
translation. Matrices are projective, and composition acts from
right to left.

\subsection{The leaf and its adjacent white faces}
Normalize the ordered base of a leaf to $(0,\infty,1,R)$.
At completeness,
\begin{equation}\label{cg:normalized}
 E=e^{2\pi\ii a},\qquad Z=\frac{1+E}{2},\qquad
 R=\frac{(1+E)^2}{4E}=\cos^2(\pi a).
\end{equation}
The apices in the two charts are $Z$ and $Z/E=\overline Z$.
Their white or completion maps are
\begin{equation}\label{cg:white}
\begin{aligned}
 W_0&=\begin{pmatrix}1&0\\0&E\end{pmatrix},&
 W_1&=\begin{pmatrix}1&E-1\\0&E\end{pmatrix},\\
 W_2&=\begin{pmatrix}E-R&R(1-E)\\E-1&1-ER\end{pmatrix},&
 W_3&=\begin{pmatrix}R&0\\1-E&ER\end{pmatrix}.
\end{aligned}
\end{equation}
The elementary pyramid has actual white faces only at $W_0,W_2$.
The expressions $W_1,W_3$ describe the maps supplied by the
adjacent prism at its cap edges. This is proved below before
using them in a group word.

The pure within-copy shaded map and the odd cross-copy maps are
\begin{equation}\label{cg:shaded}
\begin{aligned}
 S&=\begin{pmatrix}0&-ER\\1&-1-E\end{pmatrix},&
 A&=\begin{pmatrix}R&-R\\1&ER-1-E\end{pmatrix},&
 B&=\begin{pmatrix}R/E-1-1/E&R\\-1&R\end{pmatrix}.
\end{aligned}
\end{equation}
Their vertex assignments are
\[
 S:(Z,\infty,1)\mapsto(Z,0,R),\quad
 A:(Z,\infty,1)\mapsto(Z/E,R,0),\quad
 B:(Z,R,0)\mapsto(Z/E,\infty,1).
\]
Three ordered vertices determine the maps; substitution of
$Z^2-(1+E)Z+ER=0$ verifies their apex values. At the complete
parameters all these matrices are invertible.

\subsection{Compatibility across an entire prism}
Use independent phase variables $\alpha,\beta,\gamma$ for
$e^{2\pi\ii p},e^{2\pi\ii r},e^{2\pi\ii q}$.
The world coordinates are
\begin{equation}\label{cg:phase-coordinates}
\begin{split}
 v_0&=1,\quad v_1=\alpha,\quad
 v_2=\alpha\beta,\quad v_3=\alpha\beta\gamma,\\
 P&=1+\alpha-
 \frac{\alpha[\alpha\beta(1+\gamma)-1-\alpha]}
             {\alpha^2\beta^2\gamma-\alpha}.
\end{split}
\end{equation}
Indeed the chord lines have equations
$P+\alpha\overline P=1+\alpha$ and
$P+\alpha^2\beta^2\gamma\overline P=\alpha\beta(1+\gamma)$,
which give \eqref{cg:phase-coordinates}.
For an ordered quadrilateral $(a,b,c,d)$ use
\[
 M_{a,b,c}(z)=\frac{(z-a)(c-b)}{(z-b)(c-a)},\qquad
 R=M_{a,b,c}(d).
\]
Its world white map is $\overline M^{-1}W_j(E,R)M$.
Here the bar reciprocates $\alpha,\beta,\gamma$; on the complete
unit-circle chamber it is complex conjugation.
For the lower, upper and output quadrilaterals in
\eqref{cg:stack-maps} and \eqref{cg:prism-faces}, their wire
phases are $\beta\gamma,\alpha\beta,\alpha\beta\gamma$.
In these common world coordinates the entire maps obey
\begin{equation}\label{cg:nine}
\begin{aligned}
 W_{F,1}=W_{G,1}=W_{o,1}&=W_{\rm top},&
 W_{F,3}=W_{G,3}=W_{o,3}&=W_{\rm bottom},\\
 W_{F,2}&=W_{G,0},&
 W_{o,0}=W_{F,0},\qquad W_{o,2}&=W_{G,2}.
\end{aligned}
\end{equation}
The subscript $o$ means output. The top and bottom maps take
$(P,v_1,v_2)$ and $(\infty,v_0,v_3)$ to their reflected triples.
To verify \eqref{cg:nine}, insert \eqref{cg:phase-coordinates}
into $M$ and \eqref{cg:white} and cancel the projective scalar.
Each side takes the same three distinct indicated vertices to
their reflected partners. This proves equality of whole maps,
not just of their multipliers. The formulas also give nine
rational identities in the three independent variables: three
top, three bottom and the three remaining equalities.

Turns and closure have explicit compatibility as well.
Put $T_Q(z)=(R-1)/(z-1)$ and $C_Q(z)=(z-R)/(z-1)$.
Direct multiplication gives, with indices modulo four,
\begin{equation}\label{cg:turn-closing}
\begin{split}
 T_QW_{j+1}(E,R)&=W_j(E^{-1},1-R)T_Q,\\
 C_QW_j(E^{-1},R)&=W_{2-j}(E,R)C_Q.
\end{split}
\end{equation}
Induction proves that all triangular pieces of a final white face
have the same entire map. Every elementary completion expression
$W_1,W_3$ is supplied by its parent prism. At the
completed/completed root the closing seams merge white faces.
Thus each $W_j$ used at a leaf is an actual white-face map of the
closed cut. Matching $\kap$ alone would not establish this.

\subsection{The full edge words}
Write $\overline S=S(E^{-1},R)$. The pure identities are
\begin{equation}\label{cg:pure-returns}
 W_0S=\overline S W_0,\qquad
 W_2S=\overline S W_2,\qquad
 W_3S=\overline S W_1 .
\end{equation}
For an odd leaf they are
\begin{equation}\label{cg:odd-returns}
 W_0^{-1}BW_2^{-1}A=I,\quad
 W_2^{-1}BW_0^{-1}A=I,\quad
 W_1^{-1}BW_3^{-1}A=I .
\end{equation}
Direct matrix multiplication proves the projective equalities:
the first two odd products are $-R/E$ times $I$, and the last
is $(R-1)/E$ times $I$. The reverse cycles are inverses.
Every final edge is an original shaded-triangle edge, so this
list covers all edge orbits. The incidence proof of
Section~\ref{cg:assembly} gives four distinct occurrences in each
orbit. Each has angle $\pi/2$, and its full four-step return is
the identity. The metric is smooth at every finite edge, with
no axial translation hidden by the angle sum.

This is also a group check on the literal manifold. The dual
two-skeleton of the original cut balls has a generator for each
paired face and a relation for each edge cycle. Equations
\eqref{cg:pure-returns}--\eqref{cg:odd-returns} are precisely those
relations, with white maps already identified by \eqref{cg:nine}.
Thus the construction respects the original group presentation,
rather than a presentation inferred from unrelated shape equations.
